% Source-derived chapter 02: The motivic invariants
% Original source: motivic-invariants-birational-maps
\chapter{The motivic invariants}
\label{motivic-invariants-birational-maps-chapter-02}
\source{motivic-invariants-birational-maps}

\begin{remark}[Source section]
\label{motivic-invariants-birational-maps-chapter-02-source}
\source{motivic-invariants-birational-maps}
\source{motivic-invariants-birational-maps}
This chapter reproduces the corresponding section of the retrieved
LaTeX source, including its definitions, statements, examples, and
proofs. It has not yet been translated into Lean.
\notready
\end{remark}

% The source section title was: The motivic invariants
\label{sec:invariants}
\source{motivic-invariants-birational-maps}



\ssec{Invariant $c(\phi)$ with values in $\Z[\Bir_{n-1}/\k]$}
\hfill

Let $\k$ be a field and
let $X$ and $Y$ be two varieties over $\k$.
A birational map $\phi : X \dto Y$
is a morphism $U \to Y$ over $\k$
defined on a Zariski dense open subscheme
$U \subset X$
which is isomorphic onto its image.
Two birational maps $\phi,\psi : X \dto Y$
are identical
if there exists an open 
$U \subset X$ such that $\phi_{|U} = \psi_{|U}$
as morphisms.
We can compose birational maps and each
birational map $\phi$ has an inverse $\phi^{-1}$,
thereby defining a
groupoid $\ul{\Bir}/\k$ of birational classes
consisting of $\k$-varieties as objects
and birational maps as morphisms.

We say that $X$ and $Y$ are birational if there is a birational
map between them. We say that $X$ and $Y$ are stably
birational if $X \times \P_\k^k$ and $Y \times \P_\k^\ell$ are birational
for some $k, \ell \ge 0$.
The set of birational equivalence classes of 
$n$-dimensional
$\k$-varieties
is denoted by $\Bir_n/\k$.
The set of birational maps 
from $X$ to $Y$ is denoted by $\Bir(X,Y)$.

Let $\phi: X \dto Y$ be a birational map between
$n$-dimensional $\k$-varieties.
We say that $\phi$ is an isomorphism at a point 
$x \in X$ if there exists a Zariski open $U \subset X$
containing $x$ such that
$\phi_{|U}$ is an open embedding into $Y$.
Here, the point $x \in X$ is not required to be closed.
We define the exceptional set as
$$\Ex(\phi) \cnec \Set{x \in X | \phi \text{ is not an isomorphism at } x} \subset X.$$
According to the definition, $\Ex(\phi)$ is a closed
subset
containing the indeterminacy locus of $\phi$.

By definition, $\phi$ establishes
an isomorphism
\begin{equation}\label{eq:iso-ex}
\source{motivic-invariants-birational-maps}
X \setminus \Ex(\phi) \simeq Y \setminus \Ex(\phi^{-1})
\end{equation}


\begin{Def}
For every birational map $\phi: X \dto Y$ of $\k$-varieties,
we set
\begin{equation}\label{eq:def-c}
\source{motivic-invariants-birational-maps}
c(\phi) \ \ \cnec 
\sum_{\substack{y \in \Ex(\phi^{-1}) \\ \codim_Y y = 1}} [\k(y)]
 \ \ -  \sum_{\substack{x \in \Ex(\phi) \\ \codim_X x = 1}} [\k(x)] \ \ \in \ \ \Z[\Bir_{n-1}/\k]
 \end{equation}
where $\Z[\Bir_{n-1}/\k]$ 
is the free abelian group generated by $\Bir_{n-1}/\k$,
and where the function fields $\k(y)$, $\k(x)$ are identified
with the corresponding birational classes.
\end{Def}


Since both $\Ex(\phi) \subset X$ 
and $\Ex(\phi^{-1}) \subset Y$ are Zariski closed,
the sums in the definition of $c(\phi)$ are finite.
The notation $c$ stands for \emph{centers}
of the exceptional divisors \cite[Definition 2.24]{KollarMori}, 
as they are
encoded in~\eqref{eq:def-c} 
in a certain way (see e.g. Proposition~\ref{pro-cWF}).

A crucial property
of $c$ is that it defines a homomorphism
from the groupoid $\ul{\Bir}/\k$
of birational classes to
$\Z[\Bir/\k]$.

\begin{lem}
\source{motivic-invariants-birational-maps}
\notready\label{lem-homc}
\source{motivic-invariants-birational-maps}
Given birational maps $\phi : X \dto Y$ and $\psi : Y \dto Z$,
we have
$$c(\psi \circ \phi) = c(\phi) + c(\psi).$$
\end{lem}

Lemma~\ref{lem-homc} has a natural proof
based on motivic cut-and-paste,
and we postpone the proof until Subsection 
\ref{ssec-leftK},
see Theorem \ref{thm-Defc}.


Note that we work over arbitrary fields.
In particular,
the definition of $c$ and the proof of Lemma~\ref{lem-homc}
do not rely on 
any version of the resolution of singularities.
Over fields of characteristic zero 
or for surfaces over arbitrary fields
we understand birational maps better, 
thanks to
Hironaka's resolution of singularities (see e.g.~\cite{KollarRes}) 
and
Weak Factorization~\cite{WF},~\cite[Lemma 54.17.2]{stacks-project}.
The latter asserts that
every birational map $\phi : X \dto Y$
between smooth complete varieties over $\k$
admits a factorization
\begin{equation}\label{eq:wf}
\source{motivic-invariants-birational-maps}
\xymatrix{
 & Y_1 \ar[dr]^{q_1} \ar[dl]_{p_1} & & Y_2 \ar[dl]_{p_2} & \cdots & Y_{m-1} \ar[dr]^{q_{m-1}} & & Y_m  \ar[dr]^{q_m} \ar[dl]_{p_m}  \\
X & & X_1 & & \cdots & & X_{m-1} & & Y \\
}
\end{equation}
where each $p_i$ (resp. $q_i$) is a sequence of blow ups
(resp. blow downs)
along smooth irreducible closed subvarieties
of codimension $\ge 2$ or an isomorphism
(which we understand
as a blow up with empty center).

\begin{proposition}
\source{motivic-invariants-birational-maps}
\notready\label{pro-cWF}
\source{motivic-invariants-birational-maps}
If a birational map $\phi: X \dto Y$ between $n$-dimensional
smooth complete $\k$-varieties 
admits
a factorization \eqref{eq:wf},
then 
\begin{equation}\label{eqn-cWF}
\source{motivic-invariants-birational-maps}
c(\phi) = 
\sum_{Z \in \cU} [\P^{n - \dim(Z) - 1} \times Z]
  -  \sum_{T \in \cD} [\P^{n - \dim(T) - 1} \times T]  \in  \Z[\Bir_{n-1}/\k]
\end{equation}
where $\cU$ \textup(resp. $\cD$\textup)
is the set of blow up centers 
of $p_i$ \textup(resp. $q_i$\textup).

In particular, the alternating sum in~\eqref{eqn-cWF} is well-defined, namely it is independent of the
choice of weak factorization of $\phi$.

\end{proposition}

In \eqref{eqn-cWF} we use the convention that $[\emptyset] = 0$.

\begin{proof}

Proposition~\ref{pro-cWF} follows easily from
Lemma \ref{lem-homc} and the definition of $c$,
because the exceptional divisors of the $p_i$'s are of the form
$\P^{n-\dim(Z)-1} \times Z$, for $Z \in \cU$
and similarly for the $q_i$'s.
\end{proof}

We can relate $c(\phi)$ to the Galois action
on the exceptional
divisors.
Let $G_\k = \Gal(\k^{\sep}/\k)$
and let $\Kt_0(\Rep(G_\k,\Q_\ell))$ be the Grothendieck
ring of continuous finite-dimensional 
$\ell$-adic representations of $G_\k$; here
$\ell$ is a fixed prime number different from $\chr(\k)$. 
Define the group homomorphism
\begin{equation}\label{eq:sigma}
\source{motivic-invariants-birational-maps}
\gs: \Z[\Bir/\k] \to \Kt_0(\Rep(G_\k,\Q_\ell))
\end{equation}
by sending the class of a $\k$-irreducible variety 
$X$ to the permutation representation
on the irreducible components of $X_{\k^{\sep}}$
with coefficients in $\Q_\ell$.

For a smooth projective variety $X/\k$, 
we consider the N\'eron-Severi group with $\Q_\ell$-coefficients
$N(X) := \NS(X_{\k^{\sep}}) \otimes \Q_\ell$.
Then $N(X)$ is a finite-dimensional $\Q_\ell$-vector
space \cite[Theorem II.4.5]{Kollarrat} with a
continuous 
Galois action.


\begin{proposition}
\source{motivic-invariants-birational-maps}
\notready\label{prop:Picnb}
\source{motivic-invariants-birational-maps}
Let $\k$ be a field 
of characteristic zero.
For every birational map $\phi : X \dto Y$ 
between smooth projective 
varieties,
we have
$$\gs(c(\phi)) = [N(Y)] - [N(X)].$$
In particular, if $\phi \in \Bir(X)$,
then $\sigma(c(\phi)) = 0$.
\end{proposition}
\begin{proof}

By the Weak Factorization Theorem, 
the statement is reduced to the case where
$\phi: X \to Y$ is a single blow up 
with a smooth
$\k$-connected center $Z$ of codimension $\ge 2$.
In this case
the result follows
from the isomorphism of Galois representations
$N(X) \simeq N(Y) \oplus \Q_\ell[E]$,
where $\Q_\ell[E]$ stands for the permutation
representation on the geometric irreducible components
of $E$, because $\sigma(c(\phi)) = \sigma(-[E]) = -[\Q_\ell[E]]$.
\end{proof}


In the rest of this section, we 
first explain for which kinds of varieties
$c$ is identically zero,
then enhance $c(\phi)$ to 
an invariant $\ti{c}(\phi)$
which takes values in the
Grothendieck groups of varieties truncated
by the dimension.
We will explain how these invariants control the structure
of the Grothendieck ring
viewed as the 
colimit of its truncations.
The reader who is interested in examples and geometric applications
can skip these parts and go directly
to Section~\ref{sec:geometric-constructions}.


\subsection{Vanishing results for $c(\phi)$}
\hfill

Given a birational automorphism 
$\phi: X \dto X$ of variety $X$,
there are situations where $c(\phi) = 0$ is automatic.
First of all, by definition
this is always the case when $\phi$
is an isomorphism, or just a pseudo-isomorphism (isomorphism in codimension one).
This applies 
to smooth projective curves
and to all
smooth projective varieties $X$
with $K_X$ is nef (e.g. Calabi--Yau varieties),
by
Theorem~\ref{thm:BurtCY}.

In dimension two,
the following vanishing result was proven in~\cite{LSZ20}.

\begin{theorem}
\source{motivic-invariants-birational-maps}
\notready\cite[Theorem 3.4]{LSZ20}
\label{thm:surfaces-main}
\source{motivic-invariants-birational-maps}
If $\k$ is a perfect field and
$X/\k$ is a surface,
then $c(\Bir(X)) = 0$.
\end{theorem}

In dimension $3$,
the vanishing $c(\Bir(X)) = 0$ still holds in certain cases,
mostly depending on the base field $\k$.
The proof relies on the existence
of the principally
polarized intermediate Jacobian defined for
rationally connected threefolds over non-algebraically
closed
fields \cite{AchterCMVial}, \cite{AchterCMVial-2}, \cite{BenoistWittenbergPerf}.

\begin{proposition}
\source{motivic-invariants-birational-maps}
\notready\label{prop:ratconn3-fold}
\source{motivic-invariants-birational-maps}
Let $\k$ be a field of characteristic zero,
and $X$ be a smooth projective
rationally connected threefold.
The image $c(\Bir(X))$ is contained
in the subgroup of $\Z[\Bir_{2}/\k]$ generated
by classes of the form 
$[\P^1 \times C]$ where $C$ 
is an irreducible smooth curve
whose geometric
irreducible components have genus zero or one.
If $\k$ is algebraically closed, then 
$c(\Bir(X)) = 0$.
\end{proposition}

See Theorem \ref{thm:ell-curves} for genus one
curves appearing in the image of $c(\Bir(\P^3))$
over non-algebraically closed fields.




\begin{proof}
By the
Weak Factorization theorem
and Proposition~\ref{pro-cWF},
the image $c(\Bir(X))$ is contained
in the subgroup of $\Z[\Bir_2/\k]$
generated by classes $[\P^1 \times C]$
for smooth irreducible curves $C$
(if the blow up center is a closed point $Z$, the
corresponding term is $[\P^2 \times Z]
= [\P^1 \times \P^1_Z])$.

Now consider the assignment
\begin{equation}\label{eq:j}
\source{motivic-invariants-birational-maps}
[X] \mapsto \left\{\begin{array}{cl}
     \; \Jac(C), & \text{$X$ is birational to a 
     ruled surface over $C$} \\
     \; 0, & \text{otherwise} \\
\end{array}\right.
\end{equation}
for every birational class $[X] \in \Bir_2/\k$ of irreducible surfaces, where
$C$ is taken to be an irreducible smooth projective curve,
and $\Jac(C)$ is the Jacobian of $C$.
This defines a
group homomorphism 
$$j: \Z[\Bir_2/\k] \to 
\Kt_0(\PPAV/\k),$$
where $\Kt_0(\PPAV/\k)$ is the Grothendieck 
group of principally
polarized abelian varieties.


Let us prove that
for any $\phi \in \Bir(X,Y)$
where $X$ and $Y$ are
smooth projective rationally connected threefolds, we have
\begin{equation}\label{eq:j-c-jac}
\source{motivic-invariants-birational-maps}
j(c(\phi)) = [\Jac(Y)] - [\Jac(X)],
\end{equation}
where $\Jac(X)$ is the intermediate Jacobian of $X$~\cite{BenoistWittenbergPerf}.
Once again by Weak Factorization
and Lemma~\ref{lem-homc},
it suffices to check~\eqref{eq:j-c-jac}
for a single blow up along an irreducible center. 
This follows from
the blow up formulas
\cite[Lemma 2.10]{BenoistWittenbergPerf}
that $\Jac(\Bl_C(X)) \simeq \Jac(X) \times \Jac(C)$ 
when $C \subset X$ is a smooth projective curve,
and $\Jac(\Bl_P(X)) \simeq \Jac(X)$ for a point $P \in X$.
In particular, for any $\phi \in \Bir(X)$ we have
$$c(\phi) \in \Ker(j) \cap ([\P^1] \cdot \Z[\Bir_1/\k]).$$

For every irreducible smooth projective curve $C$ over $\k$,
the Jacobian $\Jac(C)$ is indecomposable over $\k$ 
as a principally polarized abelian variety.
Indeed, since $C$ is irreducible,
the $\Gal(\ol{\k}/\k)$-action on the geometric
irreducible components $C_1,\ldots,C_r$,
and therefore on the factors
$\Jac(C_{\ol{\k}}) \simeq \Jac(C_1) \times \cdots \times \Jac(C_r)$
of the decomposition of $\Jac(C_{\ol{\k}})$ is transitive.
Since each $\Jac(C_i)$ is indecomposable,
it follows from the uniqueness of the 
decomposition of principally polarized abelian varieties into irreducible ones
(see \cite[Theorem 3.3]{jordan2016unique} or \cite{DebarrePAV})
that $\Jac(C)$ is indecomposable.

Since principally polarized abelian varieties
admit unique decompositions into indecomposable ones, 
the Grothendieck 
group of principally
polarized abelian varieties
$\Kt_0(\PPAV/\k)$ is isomorphic to the free abelian group
generated by 
indecomposable ones. 
From Serre's Torelli theorem \cite[Appendix]{SerreTorelli},
it follows that $\Jac(C)$, as a principally polarized
abelian variety determines $C$ when all geometric components
of $C$ have genus $\ge 2$.
Thus the kernel
\[
\Ker(\Z[\Bir_1/\k] \overset{[\P^1] \cdot}{\to} \Z[\Bir_2/\k] 
\overset{j}\to \Kt_0(\PPAV/\k))
\]
is contained in the subgroup generated by 
birational classes of
curves with
geometric irreducible components
of genus zero or genus one.
This 
proves the first statement of Proposition~\ref{prop:ratconn3-fold}.

Finally, if $\k$ is algebraically closed,
then the Torelli theorem holds
for genus one curves as well. Thus for any
$\phi \in \Bir(X)$, 
necessarily 
$c(\phi) = m[\P^1 \times \P^1]$ for some $m \in \Z$.
It follows from Proposition \ref{prop:Picnb}
that $m = 0$.
\end{proof}

\subsection{Truncated Grothendieck groups $\Kt_0(\Var^{\le n}/\k)$}\label{ssec-tG0} 
\source{motivic-invariants-birational-maps}
\hfill

Let $\k$ be a field.
Recall that $\Bir_{n}/\k$ (resp. $\Bir/\k$)
denotes the set of birational isomorphism
classes of $n$-dimensional varieties (resp. all varieties),
and we have
the free abelian group 
$$\Z[\Bir/\k] = \bigoplus_{n \ge 0} \Z[\Bir_n/\k].$$
Similarly, let $\Z[\StBir/\k]$ denote 
the 
free abelian group
generated by the stable birational classes.
We have a surjective
homomorphism $\Z[\Bir/\k] \to 
\Z[\StBir/\k]$
sending each birational class to the corresponding
stable birational class.
Furthermore, if
$\chr(\k) = 0$, we have the Larsen-Lunts
isomorphism \cite{LarsenLunts} 
\[\xymatrix{
\Kt_0(\Var/\k) \ar[d] & \Z[\Bir/\k] \ar[d] \\
\Kt_0(\Var/\k) / (\L) \ar[r]_{ \ \ \simeq} &  \Z[\StBir/\k],
}
\]
where $\Kt_0(\Var/\k)$ is the Grothendieck ring of varieties over $\k$ and $\L = [\A^1_\k]$.
There is no natural homomorphism
between
the groups in the upper
row, even in characteristic zero.

We write $\Kt_0(\Var^{\le n}/\k)$ for the Grothendieck
group of varieties of dimension $\le n$: the generators
of $\Kt_0(\Var^{\le n}/\k)$
are
classes $[X]_n$
of schemes of finite type $X/\k$ with $\dim(X) \le n$
and the relations
are generated by
\begin{equation}\label{eq:cut-and-paste}
\source{motivic-invariants-birational-maps}
[X]_n = [U]_n + [Z]_n
\end{equation}
for every open $U \subset X$ with closed complement $Z$.
We have a sequence of natural maps
\begin{equation}\label{eqn-iota}
\source{motivic-invariants-birational-maps}
    \cdots \to \Kt_0(\Var^{\le n-1}/\k) \xto{\iota_{n-1}} \Kt_0(\Var^{\le n}/\k) \xto{\iota_{n}} \Kt_0(\Var^{\le n+1}/\k) \to \cdots.
\end{equation}
The colimit of the direct system defined by~\eqref{eqn-iota}
is the underlying group of 
the Grothendieck ring of varieties
\[
\Kt_0(\Var/\k) = \mathrm{colim}_n \; \Kt_0(\Var^{\le n}/\k).
\]

In this and the next subsections
we will express the kernel and the cokernel of
$$\Kt_0(\Var^{\le n-1}/\k) \xto{\iota_{n-1}} \Kt_0(\Var^{\le n}/\k)$$
for each $n$, 
in terms of information contained in 
the groupoid of birational classes. 
For the cokernel, we have the exact sequence
\begin{equation}\label{eq:K0-seq}
\source{motivic-invariants-birational-maps}
\Kt_0(\Var^{\le n-1}/\k) \overset{\iota_{n-1}}\to \Kt_0(\Var^{\le n}/\k) \overset{\pi_n}\to \Z[\Bir_n/\k] \to 0,
\end{equation}
where 
we define
$\pi_n$ 
on 
the classes of 
finite type schemes
of dimension $\le n$ by
\begin{equation}
\label{eq:def-pi}
\source{motivic-invariants-birational-maps}
\pi_n([X]_n) = \left\{\begin{array}{cc}
     \; [X_1] + \cdots + [X_r], & \dim(X) = n \\
     \; 0, & \dim(X) < n \\
\end{array}\right.,
\end{equation}
where $X_1,\ldots,X_r$ are the irreducible components of $X$ of dimension $n$.
The maps $\iota_{n-1}$ are not injective in general;
in Proposition \ref{prop:ker-iota} we describe their kernels,
see also Remark \ref{rem:iota-injectivity} where we explain
which ones are injective.

The following lemma 
is the starting point
for studying the Grothendieck groups inductively.
See Lemma
\ref{lem:kernelXY} 
for a reformulation of the statement using the motivic invariant
$\ti{c}$.

\begin{lemma}
\source{motivic-invariants-birational-maps}
\notready\label{lem:kernelXU}
\source{motivic-invariants-birational-maps}
Let $\Var^n/\k$ denote the set of isomorphism
classes of $n$-dimensional $\k$-varieties.
Then the natural homomorphism 
\begin{equation}\label{eq:plusprojectionU}
\source{motivic-invariants-birational-maps}
\Kt_0(\Var^{\le n-1}/\k) \oplus \Z[\Var^n/\k] \to \Kt_0(\Var^{\le n}/\k)
\end{equation}
is surjective and its kernel
admits the following presentation:
\begin{equation}\label{eq:kernel-XU}
\source{motivic-invariants-birational-maps}
\langle \ \left( [X \setminus U]_{n-1}, 
- [X] + [U] \right) \ \rangle_{U \subset X},
\end{equation}
for all $n$-dimensional
varieties $X$
and open subsets $U \subset X$. 
\end{lemma}

Recall that by our conventions, varieties are assumed
to be irreducible, while
$\Kt_0(\Var^{\le n}/\k)$
is defined with $\k$-schemes of finite type as generators. Lemma
\ref{lem:kernelXU} shows in particular that this distinction
is unimportant, and $\Kt_0(\Var^{\le n}/\k)$
could have been defined with irreducible schemes
of finite type as generators.

\begin{proof}
It is clear from definitions that 
\eqref{eq:plusprojectionU}
is surjective and that
its kernel contains
\eqref{eq:kernel-XU},
so that we have a well-defined surjective homomorphism
\begin{equation}\label{eq:plusprojectionU-new}
\source{motivic-invariants-birational-maps}
\frac{\Kt_0(\Var^{\le n-1}/\k) \oplus \Z[\Var^n/\k]}{\langle \ \left( [X \setminus U]_{n-1}, 
- [X] + [U] \right) \ \rangle_{U \subset X}} 
\overset{\beta}\longrightarrow 
\Kt_0(\Var^{\le n}/\k).
\end{equation}
To show that \eqref{eq:plusprojectionU-new} is also
injective,
we construct the inverse homomorphism $\alpha$ of $\gb$.
First we define $\alpha(X)$ for a scheme $X/\k$
of finite type of dimension $\le n$. 
Since the class $[X]_n$ only depends on the reduced structure,
we can assume that $X$ is reduced.
We stratify $X$ into 
irreducible varieties which are locally closed in $X$ and
let $X_1, \dots, X_r$
be its $n$-dimensional strata
($r = 0$ if $\dim(X) < n$).
The element
\[
\alpha(X) := 
\left([X \setminus \bigsqcup_{i=1}^r X_i ]_{n-1}, \sum_{i=1}^r [X_i]\right) 
\]
is well-defined 
(that is, independent of the choice of stratification of $X$) 
when
considered in the left-hand side of \eqref{eq:plusprojectionU-new}.
We verify that 
$\alpha(X) = \alpha(V) + \alpha(X\setminus V)$
for every open subscheme $V \subset X$, so
that $\alpha$ defines a homomorphism
from $\Kt_0(\Var^{\le n}/\k)$.
By construction $\alpha$ and $\beta$ are mutually inverse.
\end{proof}




\subsection{Invariant $\ti{c}(\phi)$ with values in $\Kt_0(\Var^{\le n-1}/\k)$}\label{ssec-leftK}
\source{motivic-invariants-birational-maps}
\hfill

The next result implies Lemma~\ref{lem-homc}.

\begin{theorem}
\source{motivic-invariants-birational-maps}
\notready\label{thm-Defc}
\source{motivic-invariants-birational-maps}
There exists a unique assignment
\[
\ti{c}: \Bir(X,Y) \to \Kt_0(\Var^{\le{n-1}}/\k)
\]
defined on the birational maps 
between $n$-dimensional varieties over $\k$ which
satisfies the following two properties.
\begin{enumerate}
    \item If $i: U \hto X$ is the inclusion of an open subset
    then $\ti{c}(i) = [X \setminus U]_{n-1}$.
    \item For every $\phi \in \Bir(X, Y)$, $\psi \in \Bir(Y,Z)$,
    we have $\ti{c}(\psi \circ \phi) = \ti{c}(\phi) + \ti{c}(\psi)$.
\end{enumerate}
Furthermore, we have $c(\phi) = \pi_{n-1}(\ti{c}(\phi))$,
where $c(\phi)$ is defined by \eqref{eq:def-c} and
$\pi_{n-1}$ is defined by \eqref{eq:def-pi}.
\end{theorem}

\begin{proof}[Proof of Theorem~\ref{thm-Defc} and
Lemma~\ref{lem-homc}]

Let $\phi \in \Bir(X, Y)$, then there is an open subset $i: U \subset X$
such that $j = \phi|_U$ is an open embedding. We have $\phi = j \circ i^{-1}$
so that properties (1) and (2) determine $\ti{c}(\phi)$ to be
\begin{equation}\label{eq:def-ct}
\source{motivic-invariants-birational-maps}
    \ti{c}(\phi) = \ti{c}(j) - \ti{c}(i) = [Y \setminus j(U)]_{n-1} - [X \setminus U]_{n-1}
\end{equation}
This shows that $\ti{c}$ is unique. 

To prove the existence,
we first check that \eqref{eq:def-ct} is well-defined, that is $\ti{c}(\phi)$
is independent of
the choice of $U \subset X$. Let $i': U' \subset X$ be another open subset
such that $j' = \phi|_{U'}$ is an open embedding.
Passing to the intersection of 
$U$ and $U'$ we may assume that $U' \subset U$.
Using the defining relations of $\Kt_0(\Var^{\le n-1}/\k)$
we obtain that
\[
\ti{c}(j') - \ti{c}(i') = [Y \setminus j'(U')]_{n-1} - [X \setminus U']_{n-1} = 
[Y \setminus j(U)]_{n-1} - [X \setminus U]_{n-1} = 
\ti{c}(j) - \ti{c}(i)
\]
hence $\ti{c}$ is well-defined.

To check (2), let $\phi \in \Bir(X, Y)$ and $\psi \in \Bir(Y,Z)$,
and choose $U \subset X$ such that both $j = \phi|_U$ and 
$j'= \psi|_{j(U)}$ are open embeddings. Then
$$\ti{c}(\psi) + \ti{c}(\phi)
= [Z \setminus j'(j(U))]_{n-1} - [Y \setminus j(U)]_{n-1}  +
[Y \setminus j(U)]_{n-1} - [X \setminus U]_{n-1} = \ti{c}(\psi 
\circ \phi).$$

For the final claim, by \eqref{eq:iso-ex} we can factorize $\phi$ through
the inclusions
\[
X \hlto X \setminus \Ex(\phi) \simeq Y \setminus \Ex(\phi^{-1}) \hto Y,
\]
so that by (1) and (2) we have
$\ti{c}(\phi) = [\Ex(\phi^{-1})]_{n-1} - [\Ex(\phi)]_{n-1}$.
It follows from the definitions that $\pi_{n-1}(\ti{c}(\phi)) = c(\phi)$. 
\end{proof}

\subsection{Invariant $\ti{c}$ and
the structure of the Grothendieck rings}
\hfill

The material of this section makes explicit
the first page of
the spectral sequences of
\cite[Section 3]{Zakharevich}
and generalizes \cite[\S 3.2]{LSZ20}.

When we put
$Y = X$ in Theorem 
\ref{thm-Defc},
we obtain a homomorphism
\begin{equation}\label{wtc-BirX}
\source{motivic-invariants-birational-maps}
\ti{c}|_{\Bir(X)} : \Bir(X) \to \Kt_0(\Var^{\le{n-1}}/\k),
\end{equation}
which, as we will see, 
is in general nonzero.
We note that it is crucial that we consider
$\Kt_0(\Var^{\le n-1}/\k)$, not the full Grothendieck
ring as the target of $\ti{c}|_{\Bir(X)}$.
Indeed, consider the homomorphism
\begin{equation}\label{eq:iota}
\source{motivic-invariants-birational-maps}
\Kt_0(\Var^{\le n-1}/\k) \overset{\iota_{n-1}}\to \Kt_0(\Var^{\le n}/\k).
\end{equation}
The following lemma shows in particular
that $\iota_{n-1} \circ \ti{c}|_{\Bir(X)} = 0$.

\begin{lemma}
\source{motivic-invariants-birational-maps}
\notready\label{lem:kernelXY}
\source{motivic-invariants-birational-maps}
The kernel of \eqref{eq:plusprojectionU}
admits the following presentation:
\begin{equation}\label{eq:kernel-XY}
\source{motivic-invariants-birational-maps}
\langle \ \left( \ti{c}(\phi), 
-[Y] + [X] \right) \ \rangle_{\phi \in \Bir(X,Y)}
\end{equation}
for all birational isomorphisms $\phi$
between all irreducible $n$-dimensional
varieties $X$, $Y$.
\end{lemma}
\begin{proof}
The subgroup \eqref{eq:kernel-XU} can be rewritten as
\begin{equation}\label{eq:kernel-UX}
\source{motivic-invariants-birational-maps}
\langle \ \left( \ti{c}(j:U \hto X), 
- [X] + [U] \right) \ \rangle_{U \subset X}
\end{equation}
where 
$U \subset X$ runs over 
open subsets of all irreducible $n$-dimensional
varieties $X$. 
It remains to show that the subgroups \eqref{eq:kernel-UX}
and \eqref{eq:kernel-XY} coincide.
It is clear that \eqref{eq:kernel-UX} is a subgroup of
\eqref{eq:kernel-XY}. Conversely, every element of
\eqref{eq:kernel-XY} can be written as a combination
of elements
from \eqref{eq:kernel-UX} after decomposing
$\phi$ as a composition of open embeddings and their inverses.
\end{proof}


\begin{proposition}
\source{motivic-invariants-birational-maps}
\notready\label{prop:ker-iota}
\source{motivic-invariants-birational-maps}
We have
\begin{equation}\label{eq:ker-iota}
\source{motivic-invariants-birational-maps}
\Ker(\iota_{n-1}) = \sum_{X \in \Bir_n/\k} \ti{c}(\Bir(X))
\end{equation}
so that there is an exact sequence
\begin{equation}\label{eq:seq-K0}
\source{motivic-invariants-birational-maps}
0 \to \Image(\ti{c})_n \to \Kt_0(\Var^{\le n-1}/\k) \overset{\iota_{n-1}}\to \Kt_0(\Var^{\le n}/\k) \overset{\pi_n}\to \Z[\Bir_n/\k] \to 0,
\end{equation}
where we write $\Image(\ti{c})_n$ for the right hand side
of \eqref{eq:ker-iota}.
\end{proposition}
\begin{proof}
$\Ker(\iota_{n-1})$ coincides with kernel of $\eqref{eq:plusprojectionU}$
intersected with $\Kt_0(\Var^{n-1}/\k)$.
Hence by Lemma \ref{lem:kernelXY} every element in
$\Ker(\iota_{n-1})$ has the form
\begin{equation}\label{eq:sum-tcphi}
\source{motivic-invariants-birational-maps}
\sum_{i=1}^r \tc(\phi_i)
\end{equation}
with $\phi_i: X_i \dto Y_i$
and $\sum_{i=1}^r ([X_i] - [Y_i]) = 0 \in \Z[\Var^n/\k]$.
Since the latter is a free abelian group on the isomorphism
classes of $n$-dimensional irreducible varieties,
there is a permutation $\sigma \in S_r$ such that
$Y_i \simeq X_{\sigma(i)}$. 
Let $i \in \{1, \dots, r\}$ and let $\ell$ be the length
of $\sigma$-orbit of $i$. Then we have a composition
\[
\psi_i: X_i \overset{\phi_i}{\dto} X_{\sigma(i)} \overset{\phi_{\sigma(i)}}{\dto} X_{\sigma^2(i)} \overset{\phi_{\sigma^2(i)}}{\dto} \dots
\overset{\phi_{\sigma^{\ell-1}(i)}}{\dto}
X_{\sigma^\ell(i)}  =  X_i. 
\]
This allows to rewrite \eqref{eq:sum-tcphi}
as a sum of $\tc(\psi_i)$, with $\psi_i \in \Bir(X_i)$,
and $i$ running over a set of representative classes
for orbits of $\sigma$ on $\{1, \dots, r\}$.
\end{proof}

\begin{remark}
\source{motivic-invariants-birational-maps}
\notready\label{rem:iota-injectivity}
\source{motivic-invariants-birational-maps}
Using Proposition \ref{prop:ker-iota}
and
Theorem \ref{thm:surfaces-main},
one can prove that 
$\iota_0$ and
$\iota_1$ are injective
for all perfect fields $\k$ \cite[Corollary 3.10]{LSZ20}.

On the other hand when $n \ge 2$, geometric constructions in the following
section 
shows in an explicit way that $\iota_n$ is not injective
over various fields $\k$;
see Theorems \ref{thm:ell-curves}, \ref{thm:K3}, 
\ref{thm:CY3}, and Theorem~\ref{thm:bir-generation}.
\end{remark}
